\begin{frame}{마르코프 성질이 성립하지 않을 때: hidden state와 재설계}
  \small
  현실은 대부분 \kb{부분관측}(partially observed)이다 ---
  에이전트가 보는 ``관측''이 진짜 상태의 완전한 투영이 아니다.
  \vskip0.3em
  \textbf{대표 예: 자율주행} --- 카메라 한 장만 관측으로 쓰면
  ``지금 얼마나 빠르게 움직이고 있었는가''라는 정보가 빠져,
  \textbf{같은 화면}에서 ``속도가 높은 상황''과 ``속도가 낮은
  상황''이 구별되지 못한다 --- 같은 $s_{\text{obs}}$ 임에도 미래
  분포가 \textbf{다르다} $\Rightarrow$ 성질이 깨진다.
  \vskip0.4em
  \begin{block}{상태 재설계의 두 가지 표준 전략}
    \kb{전략 1: 상태에 정보를 직접 넣는다} --- 최근 $k$ 프레임
    $(o_t, o_{t-1}, \ldots, o_{t-k+1})$ 을 하나의 ``확장 상태''로
    묶어 $s_t$ 로 재정의. $k$ 가 충분하면 \textbf{속도 정보가
    (프레임 차이로) 복원}되어 성질이 다시 성립. 해석은 쉽지만
    ``$k$ 를 몇으로 할지''라는 \textbf{손으로 고르는
    하이퍼파라미터}가 생긴다.
    \vskip0.3em
    \kb{전략 2: 학습된 은닉 상태로 요약한다} --- 신경망(예: RNN의
    은닉 상태 $h_t$, Week 11)이 과거 관측 시퀀스를 $h_t$ 로 압축하고
    $(o_t, h_t)$ 를 상태로 쓴다. ``요약에 충분한 정보를 담았는지''는
    수학적으로 검증할 수 없고 \textbf{경험적으로} 판단.
    이것이 \kb{부분관측 MDP(POMDP)} 라는, 이번 학기에서 직접
    다루지는 않지만 \textbf{이름은 꼭 알아 둘} 문제 영역.
  \end{block}
  \vskip0.4em
  \begin{block}{핵심}
    마르코프 성질은 ``자연이 주는 성질''이 아니라, \textbf{우리가
    상태를 어떻게 정의하는가의 설계 선택}이다. 성질이 ``깨져''
    보이면, 그건 우리의 상태 정의가 \textbf{부족했다는 신호}일
    뿐, 강화학습 자체를 포기해야 할 이유가 아니다.
  \end{block}
\end{frame}
